\section{Discussion}
\label{sec:discussion}

\subsection{Principal findings}

The study yields four principal findings. First, the effectiveness of every test-time aggregation strategy is governed by a precondition that the literature rarely states explicitly: the model must produce diverse answers at all. For the flagship models in our pool, five of six benchmarks sit in a degenerate regime in which six stochastic samples and six paraphrases return the same answer for almost every question, and in that regime self-consistency, paraphrase ensembling, and the combined ensemble reproduce the greedy baseline exactly, with the hybrid allocation differing by at most one question. This part of the degeneracy result is structural rather than empirical: when all votes coincide, every aggregation of them returns the same answer by construction. The empirical content lies in how common the coincidence is, and in its consequence that agreement-based uncertainty scores are uninformative exactly there. Second, in the high-entropy regime, which in our suite is concentrated in logical reasoning for weaker models and in selected quantitative tasks for mid-tier models, both pure strategies improve over greedy decoding when the correct answer is the modal one, while the highest-entropy cells show that majority voting can leave accuracy unchanged or even reduce it when the modal answer is wrong, in line with the plurality-vote model of Section~\ref{sec:condorcet}. Third, paraphrase agreement functions as an error signal in its own right, and its error coverage overlaps only partially with that of sample agreement: across models, the stronger source differs, so neither signal subsumes the other, and a deployment that can afford both probes obtains partially disjoint evidence, although the simple sum of the two disagreement scores does not dominate the best single signal in our data. Fourth, fragility, the fraction of questions whose answer flips under paraphrasing alone, tracks answer entropy, concentrates entirely in stem-locus paraphrasing, and appears in the same model-task cells that benefit from ensembling, which supports the interpretation that paraphrase disagreement and sampling disagreement are two observable symptoms of the same underlying instability.

\subsection{A practical allocation rule, and how it performs}

The regime structure suggests a simple procedure for practitioners with a fixed per-question budget. Probe the question first with a small number of calls, compute sample agreement and paraphrase agreement, and branch on the result. If both agreement scores equal one, additional calls of either kind are unlikely to change the answer, and the budget is better spent elsewhere; the degenerate regime is identifiable from a handful of calls at negligible cost. If sample agreement is high but paraphrase agreement is low, the model is prompt-brittle, and the remaining budget should go to paraphrases. If the reverse holds, the budget should go to resampling. If both scores are low, split the remaining budget between the two sources.

We evaluated this rule by simulation on the collected call matrices, charging the probe's four calls (two canonical samples and two paraphrases) to the adaptive condition's budget. The adaptive rule matches or exceeds every fixed six-call strategy for the models where the strategies differ at all: it returns 86.7\% at a mean cost of 4.5 calls for GPT-6-Luna, against 73.3\% for fixed SC-6, and 91.7\% at 4.7 calls for Gemini3.8-Flash, while for the degenerate flagships it terminates after the probe at a mean of 4.1 calls with the same accuracy as any full ensemble. Across all six models, an even simpler sequential rule that draws canonical samples until the first unanimity and otherwise falls back to the six-sample majority preserves SC-6 accuracy within one question everywhere at a mean budget of two to three calls. Because agreement statistics are cheap to compute online, the rules add no model training and no human labeling, and both simulations are reproducible from the released logs. The hybrid allocation, by contrast, is best understood as the fixed-budget version of the same hedge: it matches the best pure strategy when one source fails and trails it by about one question when both work, so the adaptive probe dominates it in our data at two-thirds of the cost.

\subsection{Implications for evaluation and for uncertainty estimation}

Two methodological consequences follow. For benchmark evaluation, our results add an aggregation dimension to the multi-prompt argument of Mizrahi et al. \citep{mizrahi2024state}: the same question answered under one prompt with one decode, under six paraphrases, or under twelve combined calls yields systematically different accuracies in the high-entropy regime, so reported scores should state the aggregation configuration, and comparisons across papers that aggregate differently are not like for like. The single-sample baseline S1 deserves the same reporting discipline, because greedy decoding and a single stochastic draw differ by up to ten points in the high-entropy cells of our data, in both directions. For uncertainty estimation, the degenerate regime is a warning whose scope we state precisely. Our measurements cover multiple-choice questions answered directly, where sampling-based confidence measures, including simple agreement heuristics, are uninformative when the sample distribution collapses, and our data show that this collapse is the common case for strong models on saturated tasks. Semantic-entropy methods for free-form generation \citep{kuhn2023semantic,farquhar2024detecting} face an analogous risk when the sampled paraphrase classes collapse, but our design does not test that setting, and the caution transfers by argument rather than by measurement. Paraphrase probes partially restore the signal because they perturb the input rather than the decoder, but they too inherit the collapse when the model is robust in both views, as our flagship models are on five of six benchmarks. Confidence methods should therefore be validated per regime rather than averaged over mixed test sets, and deployments should verify that the answer distribution carries variance before trusting agreement-based scores.

\subsection{Limitations}

Eight limitations bound the scope of our conclusions. First, sample size and power: fifteen questions per benchmark and completed sets of 12 to 90 questions per model support paired within-model comparisons of large effects, but the cluster-bootstrap intervals are wide, TOST does not establish equivalence for most pairs, and differences of a few points between strategies would go undetected; scaling each benchmark to hundreds of questions is the single change most likely to sharpen the conclusions. Second, provider rate limits and free-tier quotas constrained coverage unevenly across models, and the quota of Qwen3.8-Omni was exhausted before its follow-up conditions could be run; coverage follows a fixed question order rather than adaptive selection, but cross-model aggregates weight models unequally. Third, the models are commercial systems accessed through third-party endpoints: we cannot verify that the requested temperature is honored, especially for reasoning-style models, and decoding internals, post-training, and version pinning are outside our control; findings describe the accessed systems as of October 2026 and need not transfer to other versions or serving configurations, although we release per-call timestamps and token usage so the configuration is auditable. Fourth, paraphrases were generated by two models and validated by embedding similarity with a fixed threshold plus a bidirectional entailment check, but semantic equivalence of adversarial questions is not fully formalizable, and stricter thresholds or human-curated paraphrases could shift the fragility statistics. Fifth, the multiple-choice tasks use direct answering while GSM8K uses chain of thought; the aggregation comparison is internally valid within each answering mode, and the follow-up condition of Section~\ref{sec:cot} measures the mode effect directly, but absolute accuracies are not comparable across modes. Sixth, the high-entropy regime is sampled through fewer model-task cells than the degenerate regime, resting primarily on logical reasoning and selected quantitative tasks; the MMLU-Pro probe of Section~\ref{sec:mmlupro} begins to broaden this, at small scale. Seventh, error-detection metrics are computed against ensemble references that share votes with the detectors; we separate the definitional degenerate stratum and report disjoint-vote correlations, but an error reference from fully independent draws would be cleaner. Eighth, all conditions use a fixed budget of six or twelve calls with majority voting; ranked voting, weighted aggregation, or learned vote combination could change the budget-matched comparison.

\subsection{Future work}

Three extensions follow naturally. The adaptive allocation rule can be moved online, interleaving probing and allocation per question rather than simulating the split afterwards, and its threshold can be tuned against a target budget rather than fixed at unanimity. The comparison can be extended to free-form generation, where semantic clustering replaces majority voting and paraphrase ensembles interact with surface-form competition effects \citep{holtzman2021surface}, and where the degeneracy phenomenon has direct consequences for semantic-entropy hallucination detection that our multiple-choice design cannot measure. Finally, the fragility statistic can serve as a benchmark diagnostic in its own right, complementing accuracy with a measure of how strongly a model's answers depend on the surface form of the question.
